\section{GoalScope}
\label{app:goalscope}

\subsection{Setup}

\paragraph{Pivot.} For each sub-task $g_k$, we identify a \pivot position $\pi_k$: the token at which the model has committed to implementing $g_k$ but has not yet emitted any token of its implementation.\footnote{In practice $\pi_k$ is the whitespace/newline token that begins the line implementing $g_k$; \Cref{app:pivots} gives the exact tokenization.} The pivot is the natural place to ask what the model \emph{intends} to do next, before any implementation choice has been made.

\paragraph{Patchscope} \citep{ghandeharioun2024patchscope}. To read the \emph{content} of a hidden representation in natural language, we transplant it into a separate \emph{probe} prompt $\tilde{p}$ whose continuation verbalizes whatever the injected state encodes. Concretely, we replace one or more activations of $\tilde{p}$ at a chosen layer and position with a representation cached from the source, then decode $\M(\tilde{p} \mid \cdot)$ greedily. We build our intent read-out, \goalscope, on this primitive (\Cref{sec:what}).

We introduce \goalscope, a patchscope (\Cref{sec:notations:interventions}) that verbalizes the model's intent at a pivot.

Given a generation in which the model implements $G$, we run $\M$ once and cache, at the pivot $\pi_k$, the outputs of a set of \emph{goal heads} $\mathcal{H}_{\text{goal}}$ (localized as described below). We then transplant these cached activations into the probe prompt
%
\begin{quote}
\small\ttfamily
System: You are a helpful assistant proficient in programming.\\
User: What are you doing at this line of the code? Keep your answer as brief as possible.\\
Assistant:
\end{quote}
%
by replacing the corresponding head inputs to $W_O^{\ell j}$ at the probe's final token, over the layers spanned by $\mathcal{H}_{\text{goal}}$, and decode the continuation greedily. The decoded text is \goalscope's verbalization of the intent carried at $\pi_k$. \arnab{finalize: additive vs.\ replace strategy, layer range, $\alpha$; report the one used in the main results and defer the ablation to the appendix.}

\paragraph{Locating goal heads.} We identify $\mathcal{H}_{\text{goal}}$ with a differentiable head mask trained to transfer intent across contexts: activations at the pivot of a \emph{source} sub-task, patched into a \emph{destination} pivot for a different sub-task, should redirect the destination toward the source's sub-task (\Cref{sec:notations:interventions}). The heads selected by the mask are those whose pivot activations causally carry the identity of the upcoming sub-task. \arnab{cross-ref the localization method/appendix; state $|\mathcal{H}_{\text{goal}}|$ and layer range.}

\subsection{Evaluation}

\subsection{Baselines/Existing works}

\subsubsection{Function Vector}

\subsubsection{J-Lens}

\subsubsection{Activation Oracle}

\subsection{Discussion/Opinions}
\arnab{our take on the lens methods. what info in the latents we can (and probably should) expect to be verbalizable in natural language.}